% GENERATED FILE. Edit canonical lessons, then run npm run generate:books.
% canonical-source-hash: fnv1a64:a3d2da97f5c6b904
% canonical-lessons: SA-C436-abhram, SA-C436-vrstih, SA-C436-pravarsati, SA-C436-klidyati, SA-C436-picchilah

\chapter{Clouds and Showers}
\label{ch:sa-clouds-and-showers}
\hlchaptermodality{436}

\begin{chapteropening}
\textbf{By the end of this chapter:} \emph{I can say that it is clouding over and starting to rain.}

Complete the last lesson of chapter 436: I can say that it is clouding over and starting to rain.
\end{chapteropening}

\section[\texorpdfstring{\sk{अभ्रम्} (abhram)}{अभ्रम् (abhram)}]{\sk{अभ्रम्} (abhram) — a bank of cloud; a cloud}

\label{lesson:SA-C436-abhram}



\begin{quote}
Before the new one: say the Sanskrit for barter, then the Sanskrit for a storehouse.
\end{quote}

\subsection*{You'll want to know: \sk{अभ्रम्}}
\textbf{\sk{अभ्रम्}} — \emph{abhram} — \textquotedblleft{}a bank of cloud; a cloud\textquotedblright{}.

\begin{quote}
\sk{आकाशे} \sk{कृष्णम्} \sk{अभ्रम्} \sk{अस्ति।}
\end{quote}

\begin{quote}
\emph{ākāśe kṛṣṇam abhram asti.}
\end{quote}

\begin{quote}
There is a dark cloud in the sky.
\end{quote}

A neuter noun.

\begin{cousinweb}
\textbf{\sk{अभ्रम्}} (\emph{abhram}) is the word MW gives for a cloud and cloudy weather; MW reports a reading of it as \emph{ab-bhra}, \textquotedblleft{}water-bearer\textquotedblright{}; MW compares Greek \emph{ombros} and Latin \emph{imber}, \textquotedblleft{}rain\textquotedblright{} (MW).
\end{cousinweb}

\begin{grammarlens}[title={What you've built}]
One more word for getting around.
\end{grammarlens}

\subsection*{Your turn}
\begin{itemize}
\raggedright
  \item \emph{Say it:} \emph{abhram}
  \item \emph{Say it:} \emph{abhram}, once more
  \item \emph{Say it:} \emph{ākāśe kṛṣṇam abhram asti.} — There is a dark cloud in the sky.
  \item \emph{Recall:} say the Sanskrit for barter, then the Sanskrit for a storehouse, then say \emph{abhram} again
\end{itemize}

\begin{tcolorbox}[breakable,colback=teal!4,colframe=teal!35!black,title={Before you move on}]
What does \sk{अभ्रम्} mean? (\textquotedblleft{}A bank of cloud; a cloud\textquotedblright{}.) Say it once more.
\end{tcolorbox}
\section[\texorpdfstring{\sk{वृष्टिः} (vṛṣṭiḥ)}{वृष्टिः (vṛṣṭiḥ)}]{\sk{वृष्टिः} (vṛṣṭiḥ) — rainfall; a shower}

\label{lesson:SA-C436-vrstih}



\begin{quote}
Before the new one: say the Sanskrit for a storehouse, then the Sanskrit for a bank of cloud.
\end{quote}

\subsection*{You'll want to know: \sk{वृष्टिः}}
\textbf{\sk{वृष्टिः}} — \emph{vṛṣṭiḥ} — \textquotedblleft{}rainfall; a shower\textquotedblright{}.

\begin{quote}
\sk{अद्य} \sk{बहु} \sk{वृष्टिः} \sk{भविष्यति।}
\end{quote}

\begin{quote}
\emph{adya bahu vṛṣṭiḥ bhaviṣyati.}
\end{quote}

\begin{quote}
Today there will be a lot of rainfall.
\end{quote}

A feminine noun.

\begin{cousinweb}
\textbf{\sk{वृष्टिः}} (\emph{vṛṣṭiḥ}) is built on the root √vṛṣ, \textquotedblleft{}to rain\textquotedblright{}, which also gives \emph{varṣati}, \textquotedblleft{}rains\textquotedblright{}, with the ending \emph{-ti}; MW gives \emph{vṛṣṭi} as rain (MW).
\end{cousinweb}

\begin{grammarlens}[title={What you've built}]
One more word for getting around.
\end{grammarlens}

\subsection*{Your turn}
\begin{itemize}
\raggedright
  \item \emph{Say it:} \emph{vṛṣṭiḥ}
  \item \emph{Say it:} \emph{vṛṣṭiḥ}, once more
  \item \emph{Say it:} \emph{adya bahu vṛṣṭiḥ bhaviṣyati.} — Today there will be a lot of rainfall.
  \item \emph{Recall:} say the Sanskrit for a storehouse, then the Sanskrit for a bank of cloud, then say \emph{vṛṣṭiḥ} again
\end{itemize}

\begin{tcolorbox}[breakable,colback=teal!4,colframe=teal!35!black,title={Before you move on}]
What does \sk{वृष्टिः} mean? (\textquotedblleft{}Rainfall; a shower\textquotedblright{}.) Say it once more.
\end{tcolorbox}
\section[\texorpdfstring{\sk{प्रवर्षति} (pravarṣati)}{प्रवर्षति (pravarṣati)}]{\sk{प्रवर्षति} (pravarṣati) — starts to pour; begins to rain}

\label{lesson:SA-C436-pravarsati}



\begin{quote}
Before the new one: say the Sanskrit for a bank of cloud, then the Sanskrit for rainfall.
\end{quote}

\subsection*{You'll want to know: \sk{प्रवर्षति}}
\textbf{\sk{प्रवर्षति}} — \emph{pravarṣati} — \textquotedblleft{}starts to pour; begins to rain\textquotedblright{}.

\begin{quote}
\sk{सायम्} \sk{अभ्रं} \sk{सहसा} \sk{प्रवर्षति।}
\end{quote}

\begin{quote}
\emph{sāyam abhraṃ sahasā pravarṣati.}
\end{quote}

\begin{quote}
In the evening the cloud suddenly pours down rain.
\end{quote}

Past (it …): \sk{अभ्रं} \sk{सहसा} \sk{प्रवृष्टवत्}, \emph{abhraṃ sahasā pravṛṣṭavat}.

\begin{cousinweb}
\textbf{\sk{प्रवर्षति}} (\emph{pravarṣati}) is built from \emph{pra}, \textquotedblleft{}forth\textquotedblright{}, and the root √vṛṣ, \textquotedblleft{}to rain\textquotedblright{}, which also gives \emph{varṣati}, \textquotedblleft{}rains\textquotedblright{}; MW gives \emph{pra-vṛṣ} as to begin to rain, to rain, with the participle \emph{pravṛṣṭa} (MW).
\end{cousinweb}

\begin{grammarlens}[title={What you've built}]
One more word for getting around.
\end{grammarlens}

\subsection*{Your turn}
\begin{itemize}
\raggedright
  \item \emph{Say it:} \emph{pravarṣati}
  \item \emph{Say it:} \emph{pravarṣati}, once more
  \item \emph{Say it:} \emph{sāyam abhraṃ sahasā pravarṣati.} — In the evening the cloud suddenly pours down rain.
  \item \emph{Recall:} say the Sanskrit for a bank of cloud, then the Sanskrit for rainfall, then say \emph{pravarṣati} again
\end{itemize}

\begin{tcolorbox}[breakable,colback=teal!4,colframe=teal!35!black,title={Before you move on}]
What does \sk{प्रवर्षति} mean? (\textquotedblleft{}Starts to pour; begins to rain\textquotedblright{}.) Say it once more.
\end{tcolorbox}
\section[\texorpdfstring{\sk{क्लिद्यति} (klidyati)}{क्लिद्यति (klidyati)}]{\sk{क्लिद्यति} (klidyati) — gets soaked; gets wet}

\label{lesson:SA-C436-klidyati}



\begin{quote}
Before the new one: say the Sanskrit for rainfall, then the Sanskrit for starts to pour.
\end{quote}

\subsection*{You'll want to know: \sk{क्लिद्यति}}
\textbf{\sk{क्लिद्यति}} — \emph{klidyati} — \textquotedblleft{}gets soaked; gets wet\textquotedblright{}.

\begin{quote}
\sk{छत्रं} \sk{विना} \sk{बालः} \sk{वर्षायां} \sk{क्लिद्यति।}
\end{quote}

\begin{quote}
\emph{chatraṃ vinā bālaḥ varṣāyāṃ klidyati.}
\end{quote}

\begin{quote}
Without an umbrella the boy gets soaked in the rain.
\end{quote}

Past (he …): \sk{सः} \sk{वर्षायां} \sk{क्लिन्नवान्}, \emph{saḥ varṣāyāṃ klinnavān}.

\begin{cousinweb}
\textbf{\sk{क्लिद्यति}} (\emph{klidyati}) comes from the root √klid, \textquotedblleft{}to be or become wet\textquotedblright{}; MW gives the participle \emph{klinna}, \textquotedblleft{}moistened, wet\textquotedblright{} (MW).
\end{cousinweb}

\begin{grammarlens}[title={What you've built}]
One more word for getting around.
\end{grammarlens}

\subsection*{Your turn}
\begin{itemize}
\raggedright
  \item \emph{Say it:} \emph{klidyati}
  \item \emph{Say it:} \emph{klidyati}, once more
  \item \emph{Say it:} \emph{chatraṃ vinā bālaḥ varṣāyāṃ klidyati.} — Without an umbrella the boy gets soaked in the rain.
  \item \emph{Recall:} say the Sanskrit for rainfall, then the Sanskrit for starts to pour, then say \emph{klidyati} again
\end{itemize}

\begin{tcolorbox}[breakable,colback=teal!4,colframe=teal!35!black,title={Before you move on}]
What does \sk{क्लिद्यति} mean? (\textquotedblleft{}Gets soaked; gets wet\textquotedblright{}.) Say it once more.
\end{tcolorbox}
\section[\texorpdfstring{\sk{पिच्छिलः} (picchilaḥ)}{पिच्छिलः (picchilaḥ)}]{\sk{पिच्छिलः} (picchilaḥ) — slippery}

\label{lesson:SA-C436-picchilah}



\begin{quote}
Before the new one: say the Sanskrit for starts to pour, then the Sanskrit for gets soaked.
\end{quote}

\subsection*{You'll want to know: \sk{पिच्छिलः}}
\textbf{\sk{पिच्छिलः}} — \emph{picchilaḥ} — \textquotedblleft{}slippery\textquotedblright{}.

\begin{quote}
\sk{वृष्टेः} \sk{अनन्तरं} \sk{मार्गः} \sk{पिच्छिलः} \sk{भवति।}
\end{quote}

\begin{quote}
\emph{vṛṣṭeḥ anantaraṃ mārgaḥ picchilaḥ bhavati.}
\end{quote}

\begin{quote}
After the shower the road becomes slippery.
\end{quote}

\begin{cousinweb}
\textbf{\sk{पिच्छिलः}} (\emph{picchilaḥ}) is the word MW gives for slimy, slippery; MW gives no derivation, so its origin is uncertain (MW).
\end{cousinweb}

\begin{grammarlens}[title={What you've built}]
One more word for getting around.
\end{grammarlens}

\subsection*{Your turn}
\begin{itemize}
\raggedright
  \item \emph{Say it:} \emph{picchilaḥ}
  \item \emph{Say it:} \emph{picchilaḥ}, once more
  \item \emph{Say it:} \emph{vṛṣṭeḥ anantaraṃ mārgaḥ picchilaḥ bhavati.} — After the shower the road becomes slippery.
  \item \emph{Recall:} say the Sanskrit for starts to pour, then the Sanskrit for gets soaked, then say \emph{picchilaḥ} again
\end{itemize}

\begin{tcolorbox}[breakable,colback=teal!4,colframe=teal!35!black,title={Before you move on}]
What does \sk{पिच्छिलः} mean? (\textquotedblleft{}Slippery\textquotedblright{}.) Say it once more.
\end{tcolorbox}
